%% kapa: exam=1 round=22 year=2011 period=1 form=A subject=economics subject_ko=경제원론 seq=24 number=64
%% source: Q-Net 감정평가사 기출문제 (cst00309) artlSeq=5036790 file=extracted/1008/5036790/2011년 제22회 감정평가사 1차시험 기출문제/1교시.hwp
%% conversion: kordoc 4.24.51 -> tools/gen1.py (auto); review=proofread-sonnet; flags=none
\begin{question}{64}
이자율과 관련된 피셔효과(Fisher effect)의 설명으로 옳은 것은?
\begin{choices}
  \item 기대 인플레이션율이 상승하면 명목이자율은 상승한다.
  \item 피셔효과는 실질이자율에서 물가상승율을 뺀 것이다.
  \item 통화량이 증가하면 이자율은 하락한다.
  \item 소득이 증가하면 이자율은 상승한다.
  \item 통화량 증가와 이자율과는 연관이 없다.
\end{choices}
\end{question}
\answer{1}
